% TOP_PROBLEM: {"id": 1278, "title": "Dirichlet's Divisor Problem", "category_id": 1, "category": {"id": 1, "name": "number_theory", "display_name": "Number Theory", "description": "Properties of integers, prime numbers, Diophantine equations.", "slug": "number-theory", "order_index": 1, "created_at": "2026-07-31T15:26:25.670Z"}, "status": "open"}
% FIRST_POSTED: 2026-09-25
% ENTRY_KIND: statement_only
% !TeX program = lualatex
% Definition and sources only; this file is not an LLM solution attempt.
\documentclass[11pt]{article}
\usepackage[margin=25mm]{geometry}
\usepackage{fontspec}
\setmainfont{Latin Modern Roman}
\usepackage{amsmath,amssymb,mathtools}
\usepackage{unicode-math}
\setmathfont{Latin Modern Math}
\usepackage{xurl,hyperref}
\hypersetup{colorlinks=true,urlcolor=blue}
\setcounter{secnumdepth}{0}
\setlength{\parindent}{0pt}
\setlength{\parskip}{5pt}
\setlength{\emergencystretch}{3em}
\begin{document}
\section*{\texorpdfstring{Dirichlet's Divisor Problem}{Dirichlet's Divisor Problem}}\label{1278}
\subsection{Definitions and mathematical statement}
For real $x\ge2$, let $d(n)=\sum_{a\mid n}1$ and
\[
 \Delta(x)=\sum_{1\le n\le x}d(n)-x\log x-(2\gamma-1)x,
 \qquad \gamma=\lim_{N\to\infty}\left(\sum_{n=1}^N\frac1n-\log N\right).
\]
The exact selected problem is to determine
$\alpha_2=\inf\{a\in{\mathbb{R}}:\forall{\varepsilon}>0,\ \Delta(x)=O_{\varepsilon}(x^{a+{\varepsilon}})\}$.
Thus the implied constant may depend on ${\varepsilon}$ but not on $x$. Specifying this infimum is not the same as proving a bound with no ${\varepsilon}$ loss at its endpoint, nor is a mean-square bound a pointwise bound. The catalogue asks for the optimal exponent rather than explicitly adding an endpoint estimate.
\par\smallskip\textit{Formulation and concept references:} \href{https://arxiv.org/html/2506.22587v1}{[S1]}.

\subsection{Short English statement}
How small can the exponent in a uniform pointwise error estimate for the summatory divisor function be?
\subsection{Sources}
\begin{itemize}
\item[S1]Karak and Mahatab, The Piltz divisor Problem in Number Fields Using The Resonance Method (2025; journal 2026).
\url{https://arxiv.org/html/2506.22587v1}.
\textit{Formulation source.}
\item[S2]Status source for Dirichlet Divisor Problem.
\url{https://arxiv.org/html/2601.01905v2}.
\textit{Further reference; not a proof endorsement.}
\item[S3]Status source for Dirichlet Divisor Problem.
\url{https://arxiv.org/pdf/1105.6155}.
\textit{Further reference; not a proof endorsement.}
\item[S4]Status source for Dirichlet Divisor Problem.
\url{https://arxiv.org/html/2604.18624v2}.
\textit{Further reference; not a proof endorsement.}
\item[S5]Status source for Dirichlet Divisor Problem.
\url{https://arxiv.org/html/2605.21476v1}.
\textit{Further reference; not a proof endorsement.}
\end{itemize}
\end{document}
